% Partie : logique intuitionniste
% Chapitre : propositions comme types
% Section : conservation des types

\documentclass[../../../include/open-logic-section]{subfiles}

\begin{document}

\olfileid[fr]{int}{pty}{tpr}

\olsection{Conservation des types}

Les conversions par réduction et permutation étudiées
précédemment se définissent aussi directement dans
\Log{tN2i} et \Log{tN3i}.
Dans \Log{tN3i}, le type d'un terme lambda~$N$
relativement au contexte~$\Gamma$ est la formule~$!A$
pour laquelle $\Gamma \Sequent N : !A$ est prouvable.
Les réductions des termes reflètent les conversions
des preuves. Une étape de réduction conserve le type
du terme lambda dans le même contexte~$\Gamma$.
C'est la \emph{conservation des types}.

\begin{prop}
Si $N$ a le type~$!A$ relativement à $\Gamma$
et si $N \redone N'$, alors $N'$ a aussi le type~$!A$
relativement à~$\Gamma$.
\end{prop}

\begin{proof}
Établissons d'abord le lemme de substitution typée.
Si $\Gamma,x:!B \Sequent P:!C$ et
$\Gamma \Sequent Q:!B$, alors
$\Gamma \Sequent \Subst{P}{Q}{x}:!C$.
On renomme les lieurs pour éviter toute capture.
Dans \Log{tN3i}, on peut ajouter une affectation
fraîche à chaque contexte d'une dérivation. Les axiomes
et les règles restent valides après renommage des lieurs
si nécessaire. On peut donc affaiblir le typage de~$Q$.

On raisonne par récurrence sur le typage de~$P$.
Pour $x$, on utilise le typage de~$Q$. Les autres
variables gardent leur axiome. Sous une abstraction
ou dans une branche d'analyse par cas, on applique
l'hypothèse de récurrence dans le contexte étendu,
puis on réapplique la règle. Pour les autres
constructeurs, on applique l'hypothèse de récurrence
aux sous-termes dans leur contexte initial.
%
Le type final est ainsi conservé.

Ce lemme traite les réductions de fonction et de somme ;
la réduction d'un produit choisit directement l'une
des deux preuves de composante. Par exemple, l'introduction
d'un couple suivie de sa projection donne :
\[
  \bottomAlignProof
  \AxiomC{}
  \Deduce$\Gamma \fCenter N_1 : !B_1$
  \AxiomC{}
  \Deduce$\Gamma \fCenter N_2 : !B_2$
  \RightLabel{\Intro{\land}}
  \BinaryInf$\Gamma \fCenter \pair{N_1}{N_2} : !B_1 \land !B_2$
  \RightLabel{\Elim{\land}[i]}
  \UnaryInf$\Gamma \fCenter \proj{i}{\pair{N_1}{N_2}} : !B_i$
  \DisplayProof
  \quad\redone\quad
  \bottomAlignProof
  \AxiomC{}
  \Deduce$\Gamma \fCenter N_i : !B_i$
  \DisplayProof
  \]
Le terme de la preuve à droite est~$N_i$,
exactement le réduit de $\proj{i}{\pair{N_1}{N_2}}$.

Pour une \Intro{\lif} suivie d'une \Elim{\lif},
la conversion s'écrit :
\begin{multline*}
  \bottomAlignProof
  \AxiomC{$\delta_2$}
  \Deduce$x: !B_1, \Gamma \fCenter N : !B_2$
  \RightLabel{\Intro{\lif}}
  \UnaryInf$\Gamma \fCenter \lambd[\typeof{x}{!B_1}][N] : !B_1 \lif !B_2$
  \AxiomC{}
  \RightLabel{$\delta_3$}
  \Deduce$\Gamma \fCenter M : !B_1$
  \RightLabel{\Elim{\lif}}
  \BinaryInf$\Gamma \fCenter (\lambd[\typeof{x}{!B_1}][N] M) : !B_2$
  \DisplayProof
  \quad\redone\\
  \bottomAlignProof
  \AxiomC{}
  \RightLabel{$\delta_3$}
  \Deduce$\Gamma \fCenter M : !B_1$
  \RightLabel{$\Subst{\delta_2}{\delta_3}{x:!B_1}$}
  \Deduce$\Gamma \fCenter \Subst{N}{M}{x} : !B_2$
  \DisplayProof
  \end{multline*}
Le terme de la preuve à droite est
$\Subst{N}{M}{x}$, le réduit du terme de gauche.
La sous-preuve~$\delta_2$ se transforme par le lemme
de substitution : ses axiomes pour $x:!B_1$,
de la forme $\Gamma' \Sequent x:!B_1$, sont remplacés
par la preuve~$\delta_3$
de $\Gamma \Sequent M:!B_1$, avec renommage
des variables liées et maintien des hypothèses
communes. Le résultat est une preuve de
$\Gamma \Sequent \Subst{N}{M}{x}:!B_2$.

Les conversions de permutation conservent également
le type. Dans l'exemple de la projection après une
analyse par cas, chaque branche avait un type produit ;
projeter dans chaque branche donne le même type de
composante que projeter après l'analyse.
Enfin, considérons une preuve~$\delta$ et un sous-terme
redex~$M$ dont la sous-dérivation~$\delta_1$ aboutit
à $\Gamma' \Sequent M:!B$. Une réduction produit
un terme~$M'$ ; le cas à la racine fournit une
sous-dérivation~$\delta_1'$ de
$\Gamma' \Sequent M':!B$.
En remplaçant $\delta_1$ par~$\delta_1'$ puis en
réappliquant les règles qui entouraient le sous-terme,
on obtient une preuve~$\delta'$ de
$\Gamma \Sequent N':!A$.
\end{proof}

La correspondance entre les preuves et la réduction
des termes donne aussi une propriété de
\emph{progression}.
Si un terme bien typé~$N$ n'est pas en forme normale,
il contient un redex~$M$. La sous-dérivation de typage
qui aboutit à $\Gamma' \Sequent M:!B$ admet la conversion
correspondante. En remplaçant ce redex par son réduit,
on obtient un terme~$N'$ et, par conservation des types,
une preuve de $\Gamma \Sequent N':!A$.
Ainsi, un terme bien typé qui n'est pas normal peut
encore se réduire. Cette progression, jointe à la
conservation des types, signifie qu'une réduction
ne reste pas bloquée sur un terme bien typé non normal :
elle produit un autre terme bien typé de même type.
Elle n'affirme pas à elle seule que toutes les suites
de réduction terminent.

\end{document}
